% GENERATED FILE. Edit canonical lessons, then run npm run generate:books.
% canonical-source-hash: fnv1a64:d5792460a4ab9706
% canonical-lessons: PA-C327-russna, PA-C327-khijhna, PA-C327-ikalla, PA-C327-dukkh, PA-C327-hanjhu

\chapter{Hurt Feelings}
\label{ch:pa-a2-hurt-feelings}
\hlchaptermodality{327}

\begin{chapteropening}
\textbf{By the end of this chapter:} \emph{I can say when someone is upset or lonely.}

Complete the last lesson of chapter 327: I can say when someone is upset or lonely.
\end{chapteropening}

\section[\texorpdfstring{\pa{ਰੁੱਸਣਾ} (russṇā)}{ਰੁੱਸਣਾ (russṇā)}]{\pa{ਰੁੱਸਣਾ} (russṇā) — to sulk}

\label{lesson:PA-C327-russna}



\begin{quote}
Before the new one: say the Punjabi for worry, then the Punjabi for embarrassment.
\end{quote}

\subsection*{You'll want to know: \pa{ਰੁੱਸਣਾ}}
\textbf{\pa{ਰੁੱਸਣਾ}} — \emph{russṇā} — \textquotedblleft{}to sulk\textquotedblright{}.

\begin{quote}
\pa{ਮੇਰਾ} \pa{ਛੋਟਾ} \pa{ਭਰਾ} \pa{ਖਿਡੌਣੇ} \pa{ਲਈ} \pa{ਰੁੱਸਿਆ।}
\end{quote}

\begin{quote}
\emph{merā chhoṭā bharā khiḍauṇe laī russiā.}
\end{quote}

\begin{quote}
My little brother sulked over a toy.
\end{quote}

Past (she …): \pa{ਉਹ} \pa{ਮੇਰੇ} \pa{ਨਾਲ} \pa{ਰੁੱਸੀ}, \emph{oh mere nāl russī}.

\begin{cousinweb}
\textbf{\pa{ਰੁੱਸਣਾ}} (\emph{russṇā}) is inherited from Sanskrit \emph{ruṣyati}, \textquotedblleft{}is angry\textquotedblright{}; CDIAL lists Punjabi \emph{russṇā}, \textquotedblleft{}to get angry\textquotedblright{}, under it (CDIAL 10794).
\end{cousinweb}

\begin{grammarlens}[title={What you've built}]
One more word for this chapter.
\end{grammarlens}

\subsection*{Your turn}
\begin{itemize}
\raggedright
  \item \emph{Say it:} \emph{russṇā}
  \item \emph{Say it:} \emph{russṇā}, once more
  \item \emph{Say it:} \emph{merā chhoṭā bharā khiḍauṇe laī russiā.} — My little brother sulked over a toy.
  \item \emph{Recall:} say the Punjabi for worry, then the Punjabi for embarrassment, then say \emph{russṇā} again
\end{itemize}

\begin{tcolorbox}[breakable,colback=teal!4,colframe=teal!35!black,title={Before you move on}]
What does \pa{ਰੁੱਸਣਾ} mean? (\textquotedblleft{}To sulk\textquotedblright{}.) Say it once more.
\end{tcolorbox}
\section[\texorpdfstring{\pa{ਖਿਝਣਾ} (khijhṇā)}{ਖਿਝਣਾ (khijhṇā)}]{\pa{ਖਿਝਣਾ} (khijhṇā) — to get annoyed}

\label{lesson:PA-C327-khijhna}



\begin{quote}
Before the new one: say the Punjabi for embarrassment, then the Punjabi for to sulk.
\end{quote}

\subsection*{You'll want to know: \pa{ਖਿਝਣਾ}}
\textbf{\pa{ਖਿਝਣਾ}} — \emph{khijhṇā} — \textquotedblleft{}to get annoyed\textquotedblright{}.

\begin{quote}
\pa{ਰੌਲੇ} \pa{ਤੋਂ} \pa{ਦਾਦਾ} \pa{ਜੀ} \pa{ਖਿਝਦੇ} \pa{ਹਨ।}
\end{quote}

\begin{quote}
\emph{raule ton dādā jī khijhde han.}
\end{quote}

\begin{quote}
Grandfather gets annoyed at noise.
\end{quote}

Past (I …): \pa{ਮੈਂ} \pa{ਭੀੜ} \pa{ਵਿੱਚ} \pa{ਬਹੁਤ} \pa{ਖਿਝਿਆ}, \emph{maiṁ bhīṛ vich bahut khijhiā}.

\begin{cousinweb}
\textbf{\pa{ਖਿਝਣਾ}} (\emph{khijhṇā}) is inherited from Sanskrit \emph{kṣīyate}, \textquotedblleft{}is worn away, is injured\textquotedblright{}; CDIAL lists Punjabi \emph{khijhṇā}, \textquotedblleft{}to be vexed\textquotedblright{}, under it (CDIAL 3695).
\end{cousinweb}

\begin{grammarlens}[title={What you've built}]
One more word for this chapter.
\end{grammarlens}

\subsection*{Your turn}
\begin{itemize}
\raggedright
  \item \emph{Say it:} \emph{khijhṇā}
  \item \emph{Say it:} \emph{khijhṇā}, once more
  \item \emph{Say it:} \emph{raule ton dādā jī khijhde han.} — Grandfather gets annoyed at noise.
  \item \emph{Recall:} say the Punjabi for embarrassment, then the Punjabi for to sulk, then say \emph{khijhṇā} again
  \item \emph{Write it:} copy \textbf{\pa{ਖਿਝਣਾ}} once, checking the whole word
\end{itemize}

\begin{tcolorbox}[breakable,colback=teal!4,colframe=teal!35!black,title={Before you move on}]
What does \pa{ਖਿਝਣਾ} mean? (\textquotedblleft{}To get annoyed\textquotedblright{}.) Say it once more.
\end{tcolorbox}
\section[\texorpdfstring{\pa{ਇਕੱਲਾ} (ikallā)}{ਇਕੱਲਾ (ikallā)}]{\pa{ਇਕੱਲਾ} (ikallā) — alone}

\label{lesson:PA-C327-ikalla}



\begin{quote}
Before the new one: say the Punjabi for to sulk, then the Punjabi for to get annoyed.
\end{quote}

\subsection*{You'll want to know: \pa{ਇਕੱਲਾ}}
\textbf{\pa{ਇਕੱਲਾ}} — \emph{ikallā} — \textquotedblleft{}alone\textquotedblright{}.

\begin{quote}
\pa{ਛੁੱਟੀ} \pa{ਵਿੱਚ} \pa{ਮੈਂ} \pa{ਘਰ} \pa{ਵਿੱਚ} \pa{ਇਕੱਲਾ} \pa{ਸੀ।}
\end{quote}

\begin{quote}
\emph{chhuṭṭī vich maiṁ ghar vich ikallā sī.}
\end{quote}

\begin{quote}
During the holiday I was alone at home.
\end{quote}

\begin{cousinweb}
\textbf{\pa{ਇਕੱਲਾ}} (\emph{ikallā}) is inherited from a reconstructed \emph{*ekkala}, \textquotedblleft{}alone\textquotedblright{}; CDIAL lists Punjabi \emph{ikallā}, \textquotedblleft{}alone\textquotedblright{}, under it (CDIAL 2506).
\end{cousinweb}

\begin{grammarlens}[title={What you've built}]
One more word for this chapter.
\end{grammarlens}

\subsection*{Your turn}
\begin{itemize}
\raggedright
  \item \emph{Say it:} \emph{ikallā}
  \item \emph{Say it:} \emph{ikallā}, once more
  \item \emph{Say it:} \emph{chhuṭṭī vich maiṁ ghar vich ikallā sī.} — During the holiday I was alone at home.
  \item \emph{Recall:} say the Punjabi for to sulk, then the Punjabi for to get annoyed, then say \emph{ikallā} again
  \item \emph{Write it:} copy \textbf{\pa{ਇਕੱਲਾ}} once, checking the whole word
\end{itemize}

\begin{tcolorbox}[breakable,colback=teal!4,colframe=teal!35!black,title={Before you move on}]
What does \pa{ਇਕੱਲਾ} mean? (\textquotedblleft{}Alone\textquotedblright{}.) Say it once more.
\end{tcolorbox}
\section[\texorpdfstring{\pa{ਦੁੱਖ} (dukkh)}{ਦੁੱਖ (dukkh)}]{\pa{ਦੁੱਖ} (dukkh) — sorrow}

\label{lesson:PA-C327-dukkh}



\begin{quote}
Before the new one: say the Punjabi for to get annoyed, then the Punjabi for alone.
\end{quote}

\subsection*{You'll want to know: \pa{ਦੁੱਖ}}
\textbf{\pa{ਦੁੱਖ}} — \emph{dukkh} — \textquotedblleft{}sorrow\textquotedblright{}.

\begin{quote}
\pa{ਉਸ} \pa{ਖ਼ਬਰ} \pa{ਤੋਂ} \pa{ਮੈਨੂੰ} \pa{ਬਹੁਤ} \pa{ਦੁੱਖ} \pa{ਹੋਇਆ।}
\end{quote}

\begin{quote}
\emph{us khabar ton mainū̃ bahut dukkh hoiā.}
\end{quote}

\begin{quote}
That news made me very sad.
\end{quote}

A masculine noun.

\begin{cousinweb}
\textbf{\pa{ਦੁੱਖ}} (\emph{dukkh}) is inherited from Sanskrit \emph{duḥkha}, \textquotedblleft{}difficulty, pain\textquotedblright{}; CDIAL lists Punjabi \emph{dukkh}, \textquotedblleft{}sorrow, pain\textquotedblright{}, under it (CDIAL 6375).
\end{cousinweb}

\begin{grammarlens}[title={What you've built}]
One more word for this chapter.
\end{grammarlens}

\subsection*{Your turn}
\begin{itemize}
\raggedright
  \item \emph{Say it:} \emph{dukkh}
  \item \emph{Say it:} \emph{dukkh}, once more
  \item \emph{Say it:} \emph{us khabar ton mainū̃ bahut dukkh hoiā.} — That news made me very sad.
  \item \emph{Recall:} say the Punjabi for to get annoyed, then the Punjabi for alone, then say \emph{dukkh} again
  \item \emph{Write it:} copy \textbf{\pa{ਦੁੱਖ}} once, checking the whole word
\end{itemize}

\begin{tcolorbox}[breakable,colback=teal!4,colframe=teal!35!black,title={Before you move on}]
What does \pa{ਦੁੱਖ} mean? (\textquotedblleft{}Sorrow\textquotedblright{}.) Say it once more.
\end{tcolorbox}
\section[\texorpdfstring{\pa{ਹੰਝੂ} (hanjhū)}{ਹੰਝੂ (hanjhū)}]{\pa{ਹੰਝੂ} (hanjhū) — a teardrop}

\label{lesson:PA-C327-hanjhu}



\begin{quote}
Before the new one: say the Punjabi for alone, then the Punjabi for sorrow.
\end{quote}

\subsection*{You'll want to know: \pa{ਹੰਝੂ}}
\textbf{\pa{ਹੰਝੂ}} — \emph{hanjhū} — \textquotedblleft{}a teardrop\textquotedblright{}.

\begin{quote}
\pa{ਵਿਆਹ} \pa{ਵਿੱਚ} \pa{ਮੇਰੀ} \pa{ਭੈਣ} \pa{ਦੀ} \pa{ਅੱਖ} \pa{ਵਿੱਚ} \pa{ਹੰਝੂ} \pa{ਆਇਆ।}
\end{quote}

\begin{quote}
\emph{viāh vich merī bhaiṇ dī akkh vich hanjhū āiā.}
\end{quote}

\begin{quote}
At the wedding a tear came to my sister's eye.
\end{quote}

A masculine noun.

\begin{cousinweb}
\textbf{\pa{ਹੰਝੂ}} (\emph{hanjhū}) is inherited from Sanskrit \emph{aśru}, \textquotedblleft{}a tear\textquotedblright{}; CDIAL lists Punjabi \emph{hañjhū}, \textquotedblleft{}tear\textquotedblright{}, under it (CDIAL 919).
\end{cousinweb}

\begin{grammarlens}[title={What you've built}]
One more word for this chapter.
\end{grammarlens}

\subsection*{Your turn}
\begin{itemize}
\raggedright
  \item \emph{Say it:} \emph{hanjhū}
  \item \emph{Say it:} \emph{hanjhū}, once more
  \item \emph{Say it:} \emph{viāh vich merī bhaiṇ dī akkh vich hanjhū āiā.} — At the wedding a tear came to my sister's eye.
  \item \emph{Recall:} say the Punjabi for alone, then the Punjabi for sorrow, then say \emph{hanjhū} again
  \item \emph{Write it:} copy \textbf{\pa{ਹੰਝੂ}} once, checking the whole word
\end{itemize}

\begin{tcolorbox}[breakable,colback=teal!4,colframe=teal!35!black,title={Before you move on}]
What does \pa{ਹੰਝੂ} mean? (\textquotedblleft{}A teardrop\textquotedblright{}.) Say it once more.
\end{tcolorbox}
